\section*{Chapter \ref{ch:s1:running-a-multi-strategy-book} --- Running a Multi-Strategy Book}

\begin{solution}{exo:s1:running-a-multi-strategy-book:1}
The net of ten independent equal trades has a standard deviation $\sqrt{10}$ times one, against a gross of ten: netting saves $1 - 1/\sqrt{10} = 68.4\%$
(the simulation measures 68.5\%).
\end{solution}

\begin{solution}{exo:s1:running-a-multi-strategy-book:2}
$0.4 \times 6 = 2.4$ volatilities, 24\% of capital at 10\% volatility.
\end{solution}

\begin{solution}{exo:s1:running-a-multi-strategy-book:3}
Their correlation is measured mostly over calm days, when the shared factor is small noise; in a crash the shared factor dominates and all three move
with it. A low average correlation hides a large common exposure in the tail.
\end{solution}

\begin{solution}{exo:s1:running-a-multi-strategy-book:4}
34.5\%: at 5\% volatility a 5\% limit is one annual volatility, the same as a 10\% limit at 10\% volatility. What matters is the limit in the \omterm{def:s1:running-a-multi-strategy-book:multi}{pod}'s own
volatilities.
\end{solution}

\begin{solution}{exo:s1:running-a-multi-strategy-book:5}
Any \omterm{def:s1:running-a-multi-strategy-book:multi}{pod}, dead or alive, hits a tight limit soon; the dead \omterm{def:s1:running-a-multi-strategy-book:multi}{pod}'s negative drift makes it sooner (98 days at 5\%). The cost is stopping nearly every live \omterm{def:s1:running-a-multi-strategy-book:multi}{pod}
each year (91.6\% on the paths), losing their edge until they restart.
\end{solution}

\begin{solution}{exo:s1:running-a-multi-strategy-book:6}
With wealth kept above a fraction of its running maximum, the optimal risky position is proportional to the surplus above that floor: risk shrinks
smoothly as losses approach the floor and grows as the surplus is rebuilt. Halving risk at one drawdown and stopping at a second is a two-step
approximation of that line.
\end{solution}

\begin{solution}{exo:s1:running-a-multi-strategy-book:7}
With a cap of a quarter of a \omterm{def:s1:running-a-multi-strategy-book:multi}{pod}, the crash losses fall to 12.2\%, 11.1\% and 13.0\%, the largest drawdown to 28.0\% and the Sharpe ratio rises to 1.72.
In a real firm the factor map is estimated, not known: a strict cap on factors measured with error would cut \omterm{def:s1:running-a-multi-strategy-book:multi}{pods} for exposures they do not have, and
would push \omterm{def:s1:running-a-multi-strategy-book:multi}{pods} toward exposures the risk model does not see.
\end{solution}

\begin{solution}{exo:s1:running-a-multi-strategy-book:8}
Three years of mostly calm data measure the correlation of noise, not of the exposures that matter in stress; the synthetic firm's overlapping \omterm{def:s1:running-a-multi-strategy-book:multi}{pods} had
0.22 on average and lost a quarter of their capital together. Measure exposures from positions and test the firm in stress scenarios.
\end{solution}

\begin{solution}{pb:s1:running-a-multi-strategy-book:1}
\begin{enumerate}
\item A firm's combined portfolio of strategies run by teams; one team with its own capital and limits; the correlation of their P\&L, which can hide
shared exposures.
\item Ten \omterm{def:s1:running-a-multi-strategy-book:multi}{pods} at 10\% volatility, Sharpe ratios 0.5 to 1.0; \omterm{def:s1:running-a-multi-strategy-book:multi}{pods} 1--3 and 4--5 on two shared factors with three crashes; \omterm{def:s1:running-a-multi-strategy-book:multi}{pod} 10 dead after ten years.
\item By inverse trailing volatility, monthly.
\item A market-wide deleveraging in quantitative equity strategies, with unwinds on 1 and 6 August.
\item 0.22 among \omterm{def:s1:running-a-multi-strategy-book:multi}{pods} on the same factor, $-0.004$ with others; 2.4 to 2.6 volatilities lost together in each crash.
\item From positions mapped to risk-model factors, capping each factor's total exposure.
\item 1.45, 43.5\% and 31.5--33.6\% without; 1.70, 35.7\% and 20.7--22.9\% with a cap of half a \omterm{def:s1:running-a-multi-strategy-book:multi}{pod}.
\item 34.6\% when \omterm{def:s1:running-a-multi-strategy-book:multi}{pods} trade a fifth of the instruments, 68.5\% when all trade all.
\item A stop on a loss from the peak; stopping live \omterm{def:s1:running-a-multi-strategy-book:multi}{pods} and running dead ones.
\item 91.6\%, 34.5\%, 8.9\%, 2.1\% and 0.0\% on paths; 85.8\%, 38.9\%, 22.1\%, 14.7\% and 5.8\% in the firm, for 5\% to 30\%.
\item 98, 307 and 979 days at 5\%, 10\% and 15\%; never at 20\% and 30\%.
\item One sample, five candidates, and differences within the noise.
\item \textbf{Named result.} At 10\% firm volatility the largest drawdown is 43.5\% without overlap detection and 35.7\% with the shared factors capped
at half a \omterm{def:s1:running-a-multi-strategy-book:multi}{pod}; a 5\% \omterm{def:s1:running-a-multi-strategy-book:limit}{drawdown limit} stops 91.6\% of healthy \omterm{def:s1:running-a-multi-strategy-book:multi}{pods} (Sharpe ratio 0.7, 10\% volatility) within a year.
\item The optimal risky position is proportional to the surplus over a floor set as a fraction of the running maximum.
\item In the \omterm{def:s1:running-a-multi-strategy-book:multi}{pod}'s own volatilities and expected Sharpe ratio.
\item Several \omterm{def:s1:running-a-multi-strategy-book:multi}{pods} lose together, each through its own limit.
\item Positions and their factor exposures, daily P\&L, and trades for netting.
\item Book~7 measured crowding across firms; this chapter measures it inside one.
\item The \omterm{def:s1:running-a-multi-strategy-book:multi}{pods} most exposed to the factor that crashed, reduced to the cap; the others did not cause the loss.
\item A view of shared risk across \omterm{def:s1:running-a-multi-strategy-book:multi}{pods}, netting, and limits set in each \omterm{def:s1:running-a-multi-strategy-book:multi}{pod}'s own units.
\end{enumerate}
\end{solution}

\begin{solution}{iq:s1:running-a-multi-strategy-book:1}
From the \omterm{def:s1:running-a-multi-strategy-book:multi}{pod}'s expected volatility and Sharpe ratio: choose the limit so that the chance of stopping a \omterm{def:s1:running-a-multi-strategy-book:multi}{pod} with the expected edge within a year is
acceptable (a few percent), with a soft step that halves risk first; revise as the \omterm{def:s1:running-a-multi-strategy-book:multi}{pod}'s record grows.
\end{solution}

\begin{solution}{iq:s1:running-a-multi-strategy-book:2}
A shared exposure: map both books' positions to common factors (crowded styles, liquidity, a common counterparty or funding source), check whether both
hold the same names or trades under different labels, and look at other firms' known positioning.
\end{solution}

\begin{solution}{iq:s1:running-a-multi-strategy-book:3}
Compare recent performance with the distribution implied by the expected Sharpe ratio, look for a mechanism that would have killed it (crowding, costs,
regime), and accept that a Sharpe ratio drop from 0.9 to $-0.3$ takes about a year or more to detect with confidence.
\end{solution}

\begin{solution}{iq:s1:running-a-multi-strategy-book:4}
Collect \omterm{def:s1:running-a-multi-strategy-book:multi}{pods}' orders by instrument over a short window, cross opposite orders internally at a fair reference price, send only the net to the market, and
allocate fills and costs back to each \omterm{def:s1:running-a-multi-strategy-book:multi}{pod} fairly; keep an audit trail.
\end{solution}

\begin{solution}{iq:s1:running-a-multi-strategy-book:5}
Reduce risk before the limit does it for you, keeping the positions with the best edge; the firm should want risk cut smoothly as the drawdown grows,
not a \omterm{def:s1:running-a-multi-strategy-book:multi}{pod} that gambles to recover or one that dumps everything at the line.
\end{solution}

\begin{solution}{iq:s1:running-a-multi-strategy-book:6}
For $\mu = 0$, $P(D_T > L) = P(|W_T| > L) = 2\big(1 - \Phi(L/(\sigma\sqrt{T}))\big)$; for $L = 0.5\sigma$ and $T = 1$ year it is 0.617. The drawdown can
exceed $L$ at some time before $T$ and recover by $T$, so the probability that it ever exceeds $L$ within $T$ is larger; that is what a stop rule
triggers on.
\end{solution}
